\documentclass[a4j, 9pt]{ltjsarticle}

\usepackage{../mypackage}

\begin{document}

  \section{数学記号一覧}
    ここに記す記号は一般に使われる数学記号である。
    本セクション以降、本書では特に断りなしに用いるのでここで確認しておこう。
    また、各分野のより専門的な記号などに関しては各セクションで説明することにする。

    \subsection{基本演算記号}
      \subsubsection{二項関係演算子}
        
        \begin{center}
          \begin{table}[h]
            \centering
            \begin{threeparttable}
              \begin{tabular}{lcl}
              
                $\ds x + y$                                         & $\ds \cdots$  & $\ds xとy$の加法演算\\
                $\ds x - y$                                         & $\ds \cdots$  & $\ds xからy$の減法演算\\
                $\ds x \times y,　x \cdot y,　xy$                   & $\ds \cdots$  & $\ds xとy$の乗法\\
                $\ds x \div y,　\frac{x}{y},　x / y$                & $\ds \cdots$  & $\ds xとy$の除法\\
                $\ds x = y$                                         & $\ds \cdots$  & $\ds xとy$が等しい\\
                $\ds x \neq y$                                      & $\ds \cdots$  & $\ds xとy$が不一致\\
                $\ds x \approx y,　x \simeq y,　x \fallingdotseq y$ & $\ds \cdots$  & $\ds xとy$がほぼ等しい\\
                $\ds x > y$                                         & $\ds \cdots$  & $\ds xはy$より大きい\\
                $\ds x \geq y,　x \geqq y$                           & $\ds \cdots$  & $\ds xはy$と等しいかそれ以上\\
                $\ds x < y$                                         & $\ds \cdots$  & $\ds xはy$より小さい\\
                $\ds x \leq y,　x \leqq y$                         & $\ds \cdots$  & $\ds xはy$と等しいかそれ以下\\
                $\ds x \gg y$                                      & $\ds \cdots$  & $\ds xはy$に比べ非常に大きい\\
                $\ds x \ll y$                                      & $\ds \cdots$  & $\ds xはy$に比べ非常に小さい\\
                $\ds x \define y,　x \coloneqq y,　x \equiv y$      & $\ds \cdots$  & $\ds xをy$で定義する\\

              \end{tabular}
            \end{threeparttable}
          \end{table}
        \end{center}

    \subsection{論理記号}
      \subsubsection{論理演算子}

        \begin{center}
          \begin{table}[h]
            \centering
            \begin{threeparttable}
              \begin{tabular}{lcl}
              
                $\ds p \lor q$                                       & $\ds \cdots$  & $\ds pまたはq$\\
                $\ds p \land q$                                      & $\ds \cdots$  & $\ds pかつq$\\
                $\ds \lnot p,　\bar{p}$                              & $\ds \cdots$  & $\ds pでない$\\
                $\ds p \Longrightarrow q,　p \Rightarrow q$          & $\ds \cdots$  & $\ds pならばq$\\
                $\ds p \Longleftarrow q,　p \Leftarrow q$            & $\ds \cdots$  & $\ds qならばp$\\
                $\ds p \Longleftrightarrow q,　p \Leftrightarrow q$  & $\ds \cdots$  & $\ds pとq$の主張が等しい・$\ds pとq$は同値\\
                $\ds p \defineProposition q$                         & $\ds \cdots$  & $\ds p$の主張を$\ds q$で定義する\\

              \end{tabular}
            \end{threeparttable}
          \end{table}
        \end{center}

      \subsubsection{量化子}
        
        \begin{center}
          \begin{table}[h]
            \centering
            \begin{threeparttable}
              \begin{tabular}{lcl}
              
                $\ds \exists_x [p(x)],　\exists_x p(x)$    & $\ds \cdots$  & $\ds p(x)$を満たす$\ds x$が存在する\\
                $\ds \exists!_x [p(x)],　\exists!_x p(x)$  & $\ds \cdots$  & $\ds p(x)$を満たす$\ds x$が"ただ一つ"存在する\\
                $\ds \nexists_x [p(x)],　\nexists_x p(x)$  & $\ds \cdots$  & $\ds p(x)$を満たす$\ds x$が存在しない\\
                $\ds \forall_x [p(x)],　\forall_x p(x)$    & $\ds \cdots$  & 任意の$\ds x$で$\ds p(x)$が成立する\\

              \end{tabular}
            \end{threeparttable}
          \end{table}
        \end{center}

      \subsubsection{論証記法}

        \blank
        \begin{center}
          \begin{table}[h]
            \centering
            \begin{threeparttable}
              \begin{tabular}{lcl}
              
                $\ds \therefore$  & $\ds \cdots$  & すなわち\\
                $\ds \because$    & $\ds \cdots$  & なぜならば\\

              \end{tabular}
            \end{threeparttable}
          \end{table}
        \end{center}

\end{document}